\documentclass[11pt,a4paper]{article}
\usepackage[T1]{fontenc}
\usepackage[utf8]{inputenc}
\usepackage{lmodern,amsmath,amssymb,booktabs,array,xcolor,microtype,longtable}
\usepackage[a4paper,left=2.3cm,right=2.3cm,top=2.5cm,bottom=2.5cm]{geometry}
\usepackage[colorlinks=true]{hyperref}
\definecolor{ojacblue}{rgb}{0,0,0.8}
\hypersetup{linkcolor=ojacblue,citecolor=ojacblue,urlcolor=ojacblue}
\renewcommand{\baselinestretch}{1.08}
\setlength{\parindent}{0pt}
\setlength{\parskip}{4pt}
\setlength{\emergencystretch}{3em}
\newcommand{\R}{https://github.com/ralfanop/Certifications_Eulerian_polar-curve_computations}
\newcommand{\A}{\R/blob/main/certifications/manuscript15}
\newcommand{\dir}[1]{\href{\R/tree/main/#1}{\texttt{\detokenize{#1}}}}
\newcommand{\file}[1]{\href{\A/#1}{\texttt{\detokenize{#1}}}}
% A complete hyperlink is one indivisible word: when it does not fit, it moves in its entirety to the next line.
\newcommand{\WholeURL}[1]{\hfil\penalty200\hfilneg\mbox{\url{#1}}}
\begin{document}

\begin{center}
{\Large\bfseries Manuscript 15: Lean~4 and Arb certification record}\\[6pt]
{\large Cycloidal reconstruction and probabilistic rigidity of Eulerian Fourier profiles}\\[6pt]
Roberto Alfano, 22~November~2026
\end{center}

\section*{1. Scope.}
This record states, for manuscript~15, which results are kernel-checked in Lean~4~[3], which values are certified in Arb ball arithmetic~[1] and how both are rebuilt and checked again. It replaces the former Appendix~B of the article (correspondence between statements and Lean theorems). Equation, section and result numbers are those of the revision of 22~November~2026 that cites this record. All files named below are in the directory \dir{certifications/manuscript15} of the repository \WholeURL{\R}.

\section*{2. Lean~4: the libraries \texttt{PrivateApex} and \texttt{EulerianCert}.}
The formalisation comprises 34 modules, in the directory \file{LEAN}: 27 modules of the library \texttt{PrivateApex} and 7 modules of the library \texttt{EulerianCert}, the latter byte-identical to those of the same name in \dir{lean/EulerianCert}. They are compiled with the Lean~4 toolchain \texttt{v4.35.0-rc3} and Mathlib~[2] commit \texttt{0272bbe1c093d038d7257ab15306abd58245f9ff}. No module contains \texttt{sorry}, \texttt{admit}, \texttt{native\_decide} or an \texttt{axiom} declaration. The command \texttt{\#print axioms} run on each of the 508 theorems of \texttt{PrivateApex} reports \texttt{propext}, \texttt{Classical.choice} and \texttt{Quot.sound} for 505 of them, \texttt{propext} and \texttt{Quot.sound} for 2, and \texttt{propext} alone for 1. The independent checker \texttt{leanchecker} accepts all 27 modules of \texttt{PrivateApex}. The archived logs are \file{LOGS/SUMMARY.txt} (build), \file{LOGS/AXIOMS.txt} (axioms), \file{LOGS/SOURCE_AUDIT.txt} (source audit) and \file{LOGS/leanchecker.log}, with the checksums in \file{SHA256SUMS.txt}.

The archive is verified again, independently, by the workflow \href{\R/blob/main/.github/workflows/m15-certification.yml}{\texttt{m15-certification.yml}}, in \href{\R/actions/runs/45045814495}{run 45045814495} (commit \texttt{a774edaa}). The run checks the checksums and rebuilds the 34 modules from source. It prints the axioms of the 508 theorems again, requires the output to equal \texttt{LOGS/AXIOMS.txt} line by line with only the three standard axioms, and repeats the source audit. Finally it runs \texttt{leanchecker} on the 27 modules of \texttt{PrivateApex}. To reproduce locally, install the toolchain of \texttt{LEAN/lean-toolchain} with elan, enter \texttt{LEAN}, and run \texttt{lake exe cache get} and then \texttt{bash verify\_all.sh}.

\section*{3. Statements and Lean theorems.}
Each statement of the article listed below is the conclusion of the named Lean theorems (132 in all), in the modules shown. No result is assumed as an axiom: the only axioms used are the standard ones listed in Section 2.
{\small\raggedright
\begin{longtable}{@{}>{\raggedright\arraybackslash}p{0.2\textwidth}>{\raggedright\arraybackslash}p{0.46\textwidth}>{\raggedright\arraybackslash}p{0.28\textwidth}@{}}
\toprule
Statement & Lean theorems & Module(s)\\
\midrule
\endhead
\bottomrule
\endfoot

Definition~7.3 & \texttt{sOdd\_\allowbreak is\_\allowbreak slope}, \texttt{tEven\_\allowbreak is\_\allowbreak slope}, \texttt{plateau\_\allowbreak level}, \texttt{virtual\_\allowbreak apex\_\allowbreak point}, \texttt{virtualApexAngle\_\allowbreak eq} & \href{\A/LEAN/PrivateApex/ApexDefinitions.lean}{\texttt{ApexDefinitions}}\\
Theorem~7.4 & \texttt{alphaVirt\_\allowbreak unique\_\allowbreak min}, \texttt{alphaVirt\_\allowbreak twelve}, \texttt{alphaVirt\_\allowbreak three}, \texttt{alphaVirt\_\allowbreak threshold}, \texttt{slopeQ\_\allowbreak classification} & \href{\A/LEAN/PrivateApex/ApexDefinitions.lean}{\texttt{ApexDefinitions}}, \href{\A/LEAN/PrivateApex/GlobalApex.lean}{\texttt{GlobalApex}}\\
Theorem~7.5 & \texttt{laplace\_\allowbreak all\_\allowbreak orders}, \texttt{odd\_\allowbreak slope\_\allowbreak all\_\allowbreak orders}, \texttt{dC\_\allowbreak unique}, \texttt{logC\_\allowbreak table} & \href{\A/LEAN/PrivateApex/LaplaceAllOrders.lean}{\texttt{LaplaceAllOrders}}, \href{\A/LEAN/PrivateApex/SlopeAllOrders.lean}{\texttt{SlopeAllOrders}}, \href{\A/LEAN/PrivateApex/LogSincSeries.lean}{\texttt{LogSincSeries}}\\
Proposition~7.6 & \texttt{central\_\allowbreak difference\_\allowbreak add}, \texttt{slopeQ\_\allowbreak even\_\allowbreak eq}, \texttt{slopeQ\_\allowbreak even\_\allowbreak eq\_\allowbreak integral} & \href{\A/LEAN/PrivateApex/ConsecutiveSlopes.lean}{\texttt{ConsecutiveSlopes}}, \href{\A/LEAN/PrivateApex/OddSlopeTail.lean}{\texttt{OddSlopeTail}}\\
Corollary~7.7 & \texttt{even\_\allowbreak of\_\allowbreak odd\_\allowbreak hasExp}, \texttt{even\_\allowbreak slope\_\allowbreak all\_\allowbreak orders}, \texttt{eC\_\allowbreak unique}, \texttt{eDirect\_\allowbreak eq\_\allowbreak eC} & \href{\A/LEAN/PrivateApex/SlopeAllOrders.lean}{\texttt{SlopeAllOrders}}, \href{\A/LEAN/PrivateApex/EvenDirect.lean}{\texttt{EvenDirect}}\\
Table~2 & \texttt{cC\_\allowbreak values}, \texttt{dC\_\allowbreak values}, \texttt{eC\_\allowbreak values} & \href{\A/LEAN/PrivateApex/SlopeAllOrders.lean}{\texttt{SlopeAllOrders}}\\
Eqs.~(220), (222) & \texttt{slope\_\allowbreak ratio\_\allowbreak even\_\allowbreak exact}, \texttt{slope\_\allowbreak ratio\_\allowbreak even\_\allowbreak hasSum} & \href{\A/LEAN/PrivateApex/ConsecutiveRatios.lean}{\texttt{ConsecutiveRatios}}\\
Corollary~7.8 & \texttt{arctan\_\allowbreak taylor}, \texttt{odd\_\allowbreak angle\_\allowbreak all\_\allowbreak orders}, \texttt{even\_\allowbreak angle\_\allowbreak all\_\allowbreak orders}, \texttt{odd\_\allowbreak angle\_\allowbreak eighth\_\allowbreak order}, \texttt{even\_\allowbreak angle\_\allowbreak eighth\_\allowbreak order} & \href{\A/LEAN/PrivateApex/AngleAllOrders.lean}{\texttt{AngleAllOrders}}\\
Table~3 & \texttt{angQ\_\allowbreak dC\_\allowbreak table}, \texttt{angQ\_\allowbreak eC\_\allowbreak table}, \texttt{dC\_\allowbreak eq\_\allowbreak dTab}, \texttt{eC\_\allowbreak eq\_\allowbreak eTab} & \href{\A/LEAN/PrivateApex/AngleAllOrders.lean}{\texttt{AngleAllOrders}}\\
Proposition~7.9 & \texttt{slope\_\allowbreak ratio\_\allowbreak odd\_\allowbreak hasExp}, \texttt{slopeRatioOddQ\_\allowbreak values}, \texttt{deficit\_\allowbreak ratio\_\allowbreak even\_\allowbreak hasExp}, \texttt{deficit\_\allowbreak ratio\_\allowbreak odd\_\allowbreak hasExp}, \texttt{deficitRatioEvenC\_\allowbreak one}, \texttt{deficitRatioEvenC\_\allowbreak two}, \texttt{deficitRatioOddC\_\allowbreak one}, \texttt{deficitRatioOddC\_\allowbreak two} & \href{\A/LEAN/PrivateApex/ConsecutiveRatios.lean}{\texttt{ConsecutiveRatios}}, \href{\A/LEAN/PrivateApex/DeficitRatioCoefficients.lean}{\texttt{DeficitRatioCoefficients}}\\
Eqs.~(239), (240) & \texttt{deficitRatioEvenC\_\allowbreak three}, \texttt{deficitRatioEvenC\_\allowbreak four}, \texttt{deficitRatioOddC\_\allowbreak three}, \texttt{deficitRatioOddC\_\allowbreak four} & \href{\A/LEAN/PrivateApex/DeficitRatioCoefficients.lean}{\texttt{DeficitRatioCoefficients}}\\
Theorem~7.10 & \texttt{odd\_\allowbreak slope\_\allowbreak remainder}, \texttt{even\_\allowbreak slope\_\allowbreak remainder}, \texttt{explicit\_\allowbreak constants} & \href{\A/LEAN/PrivateApex/QuantSlopes.lean}{\texttt{QuantSlopes}}\\
Theorem~7.10 (continued) & \texttt{odd\_\allowbreak slope\_\allowbreak remainder\_\allowbreak large}, \texttt{even\_\allowbreak slope\_\allowbreak remainder\_\allowbreak large}, \texttt{odd\_\allowbreak slope\_\allowbreak remainder\_\allowbreak power}, \texttt{even\_\allowbreak slope\_\allowbreak remainder\_\allowbreak power} & \href{\A/LEAN/PrivateApex/QuantSlopes.lean}{\texttt{QuantSlopes}}\\
Eqs.~(243)--(245) & \texttt{sin\_\allowbreak taylor\_\allowbreak sharp}, \texttt{log\_\allowbreak hasExpQ}, \texttt{exp\_\allowbreak neg\_\allowbreak taylor}, \texttt{kap\_\allowbreak pow\_\allowbreak bound}, \texttt{laplace\_\allowbreak explicit} & \href{\A/LEAN/PrivateApex/QuantSinc.lean}{\texttt{QuantSinc}}, \href{\A/LEAN/PrivateApex/LaplaceAllOrders.lean}{\texttt{LaplaceAllOrders}}, \href{\A/LEAN/PrivateApex/QuantLaplace.lean}{\texttt{QuantLaplace}}\\
Theorem~7.16 & \texttt{cardioid\_\allowbreak sub\_\allowbreak quarter}, \texttt{cardioid\_\allowbreak eulerian}, \texttt{norm\_\allowbreak cardioid\_\allowbreak sub}, \texttt{eulerianB\_\allowbreak exp}, \texttt{norm\_\allowbreak eulerianB\_\allowbreak exp}, \texttt{cardioid\_\allowbreak eulerianB}, \texttt{rhoE\_\allowbreak sq\_\allowbreak mul}, \texttt{rhoE\_\allowbreak sq\_\allowbreak mul\_\allowbreak neg\_\allowbreak pi}, \texttt{rhoPole\_\allowbreak sq\_\allowbreak mul}, \texttt{one\_\allowbreak sub\_\allowbreak exp\_\allowbreak eq\_\allowbreak eulerianB}, \texttt{fixed\_\allowbreak point}, \texttt{multiplier} & \href{\A/LEAN/PrivateApex/CardioidIdentities.lean}{\texttt{CardioidIdentities}}\\
Proposition~7.17 & \texttt{z0\_\allowbreak sq}, \texttt{cardioid\_\allowbreak special\_\allowbreak point}, \texttt{special\_\allowbreak point\_\allowbreak unique} & \href{\A/LEAN/PrivateApex/CardioidIdentities.lean}{\texttt{CardioidIdentities}}\\
Theorem~7.19 & \texttt{charFun\_\allowbreak unifLaw}, \texttt{charFun\_\allowbreak unif01}, \texttt{charFun\_\allowbreak unif01\_\allowbreak eulerianB}, \texttt{triLaw\_\allowbreak eq}, \texttt{cardioid\_\allowbreak probabilistic}, \texttt{cardioid\_\allowbreak expectation} & \href{\A/LEAN/PrivateApex/CardioidProbability.lean}{\texttt{CardioidProbability}}\\
Theorem~7.20 & \texttt{uniform\_\allowbreak reconstructs}, \texttt{rigidity\_\allowbreak local}, \texttt{rigidity\_\allowbreak global}, \texttt{rigidity\_\allowbreak iff} & \href{\A/LEAN/PrivateApex/CardioidRigidity.lean}{\texttt{CardioidRigidity}}\\
Proposition~7.22 & \texttt{hasDerivAt\_\allowbreak cardioid}, \texttt{cardioid\_\allowbreak zero}, \texttt{cardioid\_\allowbreak pi}, \texttt{cardioid\_\allowbreak re}, \texttt{cardioid\_\allowbreak im}, \texttt{cusp\_\allowbreak quotient\_\allowbreak tendsto}, \texttt{period\_\allowbreak two\_\allowbreak factor}, \texttt{period\_\allowbreak doubling}, \texttt{cusp\_\allowbreak distance\_\allowbreak rE}, \texttt{rE\_\allowbreak pi} & \href{\A/LEAN/PrivateApex/CardioidIdentities.lean}{\texttt{CardioidIdentities}}\\
Eqs.~(274)--(278) & \texttt{z1\_\allowbreak eq\_\allowbreak two\_\allowbreak mul\_\allowbreak z0}, \texttt{z1\_\allowbreak eq\_\allowbreak neg\_\allowbreak inv\_\allowbreak z0}, \texttt{z0\_\allowbreak mul\_\allowbreak z1}, \texttt{z1\_\allowbreak sq}, \texttt{z0\_\allowbreak add\_\allowbreak z1}, \texttt{special\_\allowbreak quadratic}, \texttt{quadratic\_\allowbreak family\_\allowbreak product}, \texttt{quadratic\_\allowbreak family\_\allowbreak involution}, \texttt{quadratic\_\allowbreak family\_\allowbreak factor}, \texttt{golden\_\allowbreak quadratic}, \texttt{golden\_\allowbreak other\_\allowbreak root}, \texttt{S\_\allowbreak smul\_\allowbreak z0H}, \texttt{S\_\allowbreak smul\_\allowbreak z1H}, \texttt{modular\_\allowbreak invariant\_\allowbreak eq} & \href{\A/LEAN/PrivateApex/CardioidModular.lean}{\texttt{CardioidModular}}\\
Proposition~7.24 & \texttt{cardioidFam\_\allowbreak z0}, \texttt{cardioidFam\_\allowbreak z1}, \texttt{cardioidFam\_\allowbreak eulerian}, \texttt{cardioid\_\allowbreak homothety} & \href{\A/LEAN/PrivateApex/CardioidModular.lean}{\texttt{CardioidModular}}\\
Eqs.~(283), (284) & \texttt{norm\_\allowbreak cardioidFam\_\allowbreak z0}, \texttt{norm\_\allowbreak cardioidFam\_\allowbreak z1}, \texttt{cardioidFam\_\allowbreak pi}, \texttt{cardioidFam\_\allowbreak pi\_\allowbreak z0}, \texttt{cardioidFam\_\allowbreak pi\_\allowbreak z1}, \texttt{cardioidFam\_\allowbreak half\_\allowbreak pi\_\allowbreak z0}, \texttt{cardioidFam\_\allowbreak half\_\allowbreak pi\_\allowbreak z1}, \texttt{half\_\allowbreak pi\_\allowbreak ray} & \href{\A/LEAN/PrivateApex/CardioidModular.lean}{\texttt{CardioidModular}}\\
Eqs.~(285), (286) & \texttt{H4\_\allowbreak eq}, \texttt{H4\_\allowbreak cusp}, \texttt{H4\_\allowbreak cardioid}, \texttt{hasDerivAt\_\allowbreak H4}, \texttt{H4\_\allowbreak dist}, \texttt{H4\_\allowbreak dist\_\allowbreak sq}, \texttt{H4\_\allowbreak pullback} & \href{\A/LEAN/PrivateApex/CardioidModular.lean}{\texttt{CardioidModular}}\\
Eqs.~(285), (286) (continued) & \texttt{H4\_\allowbreak spherical\_\allowbreak dilation\_\allowbreak cusp}, \texttt{H4inf\_\allowbreak eq}, \texttt{hasDerivAt\_\allowbreak H4inf}, \texttt{H4\_\allowbreak spherical\_\allowbreak dilation\_\allowbreak infinity}, \texttt{H4\_\allowbreak not\_\allowbreak modular} & \href{\A/LEAN/PrivateApex/CardioidModular.lean}{\texttt{CardioidModular}}\\
\end{longtable}}

\section*{4. Arb certificates at 3600000 bits.}
Both certificates use FLINT/Arb~[1] through python-flint 0.9.0. The same run 45045814495 recomputed both of them and found the archived results again byte for byte: the enclosures, the exact binary endpoints and the decimal files.

\emph{Carlini--Laplace angle, equations (295)--(296).} The unique positive root $u_*$ of $u\tanh u=1$ is enclosed between exact dyadic endpoints with interval signs of opposite strict sign. Its uniqueness follows from $\tanh u+u\operatorname{sech}^2u>0$ for $u>0$. The enclosure is propagated to $\lambda=(u_*^2-1)^{1/2}$, $\alpha=2\arctan(1/\lambda)$ and to degrees, and checked against equivalent formulas in $\sinh u_*$ by overlap. The working precision is 3600000 bits and the relative accuracy is at least 3599870 bits. The printed 50-place value $112.93167025490469621261195929014436025897847010921393^\circ$ is certified under both truncation and rounding. This is a numerical enclosure, not a Lean proof of the catenoid theorem. Files: \file{ARB/M15_Eq295_FLINT_Arb_3600000bits.zip}, \file{ARB/Carlini-Laplace/certify_M15_eq295.py}, \file{ARB/Carlini-Laplace/M15_eq295_run.txt}, \file{ARB/Carlini-Laplace/M15_eq295_README.txt}.

\emph{Minimum apex angle, Theorem~7.4.} The angle $2\arctan(3628800/9473633)$ is computed at 3600000 bits and checked by overlap against the arccosine formula. Its relative accuracy is 3599998 bits in radians and 3599996 bits in degrees. Exact rational comparisons of the endpoints certify 1000001 decimal places, in radians and in degrees, under both truncation and rounding. The archive \file{ARB/M15_virtual_apex_3600000bits.zip} contains the exact binary endpoints, the four decimal files, the scripts and the logs. The original certificate refers to Theorem~7.3, which is Theorem~7.4 after the insertion of Definition~7.3.

\section*{5. Exact rational and symbolic checks.}
Four further scripts run in the same workflow:
\begin{itemize}\setlength{\itemsep}{1pt}\raggedright
\item \file{verify_virtual_apex.py} repeats the exact comparisons of Theorem~7.4 and checks the displayed 50-place decimal by rational alternating-series bounds.
\item \file{verify_all_order_jets.py} checks the all-order jets.
\item \file{verify_cusp_sixth_moment.py} checks the sixth-order cusp coefficients.
\item \file{verify_eulerian_angles_n3_17_and_limit.py} checks the Eulerian angles for $n=3,\dots,17$ and their limit.
\end{itemize}

\section*{6. Limits of this record.}
This record does not cover:
\begin{itemize}\setlength{\itemsep}{1pt}\raggedright
\item the identification of the triangular density;
\item the principal-root formulation;
\item the stability inequality;
\item the numerical decimal evaluation of the minimum angle, which is certified in Arb only (Section~4);
\item the additional period-two disc derivation;
\item the catenoidal reconstruction section;
\item the conjecture on optimal remainder constants.
\end{itemize}
A conjecture is not a verified theorem.

\section*{7. Other repository evidence.}
The following earlier workflow runs concern parts of manuscript~15. Each has its own scope, and the list does not claim that every extension of the article is certified.
\begin{itemize}\setlength{\itemsep}{1pt}\raggedright

\item \href{\R/pull/325}{PR~\#325}, successful run \href{\R/actions/runs/43724374973}{43724374973}: Lean and Arb checks.
\item \href{\R/pull/338}{PR~\#338}, successful run \href{\R/actions/runs/43921559615}{43921559615}: geometry and magnetic checks.
\item \href{\R/pull/378}{PR~\#378}, successful run \href{\R/actions/runs/44296095160}{44296095160}: Lean and Arb checks.
\item \href{\R/pull/435}{PR~\#435}, successful run \href{\R/actions/runs/44494713791}{44494713791}: terminal Taylor jets (partial scope).
\item \href{\R/pull/437}{PR~\#437}, successful run \href{\R/actions/runs/44495582023}{44495582023}: Möbius spectral sine-integral.
\item \href{\R/pull/457}{PR~\#457}, successful run \href{\R/actions/runs/44518261905}{44518261905}: exponential atom masses, $W_\infty$.
\item \href{\R/pull/464}{PR~\#464}, successful run \href{\R/actions/runs/44519321578}{44519321578}: Gaussian integrals.
\item \href{\R/pull/470}{PR~\#470}, successful run \href{\R/actions/runs/44743294532}{44743294532}: exact local $W_1/W_2$ costs, not global sums.
\item \href{\R/pull/491}{PR~\#491}, successful run \href{\R/actions/runs/44918198103}{44918198103}: inverse recurrence, not closed Lagrange--B\"urmann.
\item \href{\R/pull/497}{PR~\#497}, run \href{\R/actions/runs/45034350299}{45034350299}: failed; not counted as certified.
\end{itemize}

\section*{References.}
\begin{enumerate}\setlength{\itemsep}{1pt}
\item[{[1]}] F. Johansson, Arb: efficient arbitrary-precision midpoint-radius interval arithmetic, \emph{IEEE Trans. Comput.} \textbf{66} (2017), no.~8, 1281--92. \WholeURL{https://doi.org/10.1109/TC.2017.2690633}.
\item[{[2]}] The mathlib Community, The Lean mathematical library, in \emph{Proceedings of the 9th ACM SIGPLAN International Conference on Certified Programs and Proofs (CPP 2020)}, ACM, New York, 2020, 367--81. \WholeURL{https://doi.org/10.1145/3372885.3373824}.
\item[{[3]}] L. de Moura and S. Ullrich, The Lean~4 theorem prover and programming language, in \emph{Automated Deduction -- CADE 28}, A. Platzer and G. Sutcliffe (eds.), Lecture Notes in Computer Science \textbf{12699}, Springer, Cham, 2021, 625--35. eISBN 9783030798765. \WholeURL{https://doi.org/10.1007/978-3-030-79876-5_37}.
\end{enumerate}

\end{document}
